\scanpage{11}
Sur les groupes de Galois attachés aux groupes $p$-divisibles\hfill 125

$E_1=V_C(1)$. Soit $H$ la composante neutre du $C$-groupe algébrique $G_{\mathrm{alg}}(C)$ (du fait que $C$ est algébriquement clos, nous nous permettons d'identifier un $C$-groupe algébrique à l'ensemble de ses points rationnels). Tout revient à montrer que $H=\mathbf{GL}_E$. Or, $H$ jouit des trois propriétés suivantes:

(a) \textit{$H$ est un sous-groupe réductif connexe de $\mathbf{GL}_E$, de commutant réduit aux homothéties.}

(b) \textit{$H$ contient le groupe $\Phi$ formé des automorphismes de $E$ qui sont l'identité sur $E_0$ et une homothétie sur $E_1$.}

Cela résulte du théorème 1.

(c) \textit{Les dimensions $n_0$ et $n_1$ de $E_0$ et $E_1$ sont premières entre elles.}

Le théorème 3 est donc une conséquence du résultat suivant, qui est un pur énoncé de théorie des groupes algébriques:

\textbf{Proposition 5.} -- \textit{Tout sous-groupe algébrique $H$ de $\mathbf{GL}_E$ vérifiant les conditions (a), (b), (c) ci-dessus est égal à $\mathbf{GL}_E$.}

La démonstration comporte plusieurs étapes:

1) Soit $T$ (resp. $S$) la composante neutre du centre de $H$ (resp. son groupe des commutateurs). On a $H=T\cdot S$, $T\cap S$ est fini, et $S$ est semi-simple. Vu (a), $T$ est, soit réduit à $\{1\}$, soit égal au groupe $\mathbf G_m$ des homothéties. Dans le premier cas, $H=S$ serait contenu dans le groupe unimodulaire $\mathbf{SL}_E$, ce qui est absurde car les éléments de $\Phi$ ne sont pas tous de déterminant 1. On a donc $T=\mathbf G_m$ et il va nous suffire de prouver que $S=\mathbf{SL}_E$.

2) Soit $\Theta$ le tore de dimension 2 formé des automorphismes de $E$ qui sont égaux à une homothétie sur $E_0$ et à une autre homothétie sur $E_1$. On a $\Theta=\mathbf G_m\cdot\Phi$, d'où $\Theta\subset H$ d'après ce qui précède. Soit $\psi:C^*\to\Theta$ l'homomorphisme défini par:
\[
\begin{array}{lll}
\psi(\lambda)x=\lambda^{n_1}x&\text{si}&x\in E_0\\
\psi(\lambda)x=\lambda^{-n_0}x&\text{si}&x\in E_1.
\end{array}
\]

On a $\det\psi(\lambda)=\lambda^{n_0n_1-n_1n_0}=1$ pour tout $\lambda\in C^*$. L'image $\Psi$ de $\psi$ est donc un sous-groupe connexe de $H\cap\mathbf{SL}_E$, donc aussi un sous-groupe de $S$.

3) Montrons maintenant que $S$ est \textit{simple} (i.e. tout sous-groupe algébrique distingué de $S$ est, soit fini, soit égal à $S$). Sinon, en effet, on aurait $S=(S'\times S'')/N$, avec $S'$, $S''$ semi-simples non réduits à $\{1\}$, et $N$ sous-groupe fini du centre de $S'\times S''$. Le $(S'\times S'')$-module $E$ peut s'écrire comme produit tensoriel:
\[
E=E'\otimes E'',
\]
\margin{10}
